% Options for packages loaded elsewhere
\PassOptionsToPackage{unicode}{hyperref}
\PassOptionsToPackage{hyphens}{url}
\documentclass[
  ignorenonframetext,
]{beamer}
\newif\ifbibliography
\usepackage{pgfpages}
\setbeamertemplate{caption}[numbered]
\setbeamertemplate{caption label separator}{: }
\setbeamercolor{caption name}{fg=normal text.fg}
\beamertemplatenavigationsymbolsempty
% remove section numbering
\setbeamertemplate{part page}{
  \centering
  \begin{beamercolorbox}[sep=16pt,center]{part title}
    \usebeamerfont{part title}\insertpart\par
  \end{beamercolorbox}
}
\setbeamertemplate{section page}{
  \centering
  \begin{beamercolorbox}[sep=12pt,center]{section title}
    \usebeamerfont{section title}\insertsection\par
  \end{beamercolorbox}
}
\setbeamertemplate{subsection page}{
  \centering
  \begin{beamercolorbox}[sep=8pt,center]{subsection title}
    \usebeamerfont{subsection title}\insertsubsection\par
  \end{beamercolorbox}
}
% Prevent slide breaks in the middle of a paragraph
\widowpenalties 1 10000
\raggedbottom
\AtBeginPart{
  \frame{\partpage}
}
\AtBeginSection{
  \ifbibliography
  \else
    \frame{\sectionpage}
  \fi
}
\AtBeginSubsection{
  \frame{\subsectionpage}
}
\usepackage{iftex}
\ifPDFTeX
  \usepackage[T1]{fontenc}
  \usepackage[utf8]{inputenc}
  \usepackage{textcomp} % provide euro and other symbols
\else % if luatex or xetex
  \usepackage{unicode-math} % this also loads fontspec
  \defaultfontfeatures{Scale=MatchLowercase}
  \defaultfontfeatures[\rmfamily]{Ligatures=TeX,Scale=1}
\fi
\usepackage{lmodern}
\ifPDFTeX\else
  % xetex/luatex font selection
\fi
% Use upquote if available, for straight quotes in verbatim environments
\IfFileExists{upquote.sty}{\usepackage{upquote}}{}
\IfFileExists{microtype.sty}{% use microtype if available
  \usepackage[]{microtype}
  \UseMicrotypeSet[protrusion]{basicmath} % disable protrusion for tt fonts
}{}
\makeatletter
\@ifundefined{KOMAClassName}{% if non-KOMA class
  \IfFileExists{parskip.sty}{%
    \usepackage{parskip}
  }{% else
    \setlength{\parindent}{0pt}
    \setlength{\parskip}{6pt plus 2pt minus 1pt}}
}{% if KOMA class
  \KOMAoptions{parskip=half}}
\makeatother
\usepackage{longtable,booktabs,array}
\newcounter{none} % for unnumbered tables
\usepackage{calc} % for calculating minipage widths
\usepackage{caption}
% Make caption package work with longtable
\makeatletter
\def\fnum@table{\tablename~\thetable}
\makeatother
\setlength{\emergencystretch}{3em} % prevent overfull lines
\providecommand{\tightlist}{%
  \setlength{\itemsep}{0pt}\setlength{\parskip}{0pt}}
\usepackage{bookmark}
\IfFileExists{xurl.sty}{\usepackage{xurl}}{} % add URL line breaks if available
\urlstyle{same}
\hypersetup{
  pdftitle={Synthetic and Artificial Intelligence Implications},
  hidelinks,
  pdfcreator={LaTeX via pandoc}}

\title{\texorpdfstring{Synthetic and Artificial Intelligence
Implications}{Synthetic and Artificial Intelligence Implications}}
\author{\texorpdfstring{}{}}
\date{}

\begin{document}
\frame{\titlepage}

\section{Synthetic and Artificial Intelligence: Categorical
Foundations}\label{sec:ai-implications}

\begin{frame}{Compositional Structure in Multi-Agent AI Communication}
\protect\phantomsection\label{compositional-structure-in-multi-agent-ai-communication}
The categorical framework developed in the preceding sections---case
categories, functorial semantics, enriched structure, and diagrammatic
reasoning---addresses a core challenge in contemporary artificial
intelligence: how to endow autonomous agents with structured,
compositional communication.

Current multi-agent AI systems communicate through increasingly
standardized protocols. Google's Agent-to-Agent (A2A) protocol
{[}@google2025a2a{]}, introduced in 2025, provides a standard way for
agents to communicate regardless of underlying framework through
HTTP/JSON-RPC message passing with capability discovery. The Model
Context Protocol (MCP) {[}@anthropic2024mcp{]} standardizes how AI
agents access external tools and data sources. The Agent Communication
Protocol (ACP) {[}@acp2025protocol{]} handles standardizing messaging
formats across various users, including agents, applications, and
humans. The Agent Network Protocol (ANP) {[}@anp2025protocol{]}
introduces a three-layer architecture supporting trusted agent
interaction in distributed systems. These engineering advances solve
critical interoperability problems but lack a compositional semantics:
messages are flat JSON payloads without inherent algebraic structure.

Category theory provides exactly the missing layer. Just as the DisCoCat
framework maps grammatical derivations functorially into vector spaces
(see \autoref{sec:categorical-semantics}), a categorical communication
protocol would map agent interaction patterns functorially into
behavioral specifications. The morphisms of such a protocol category
would be message types, the objects would be agent states, and the
composition law would specify how multi-step interactions chain
together---exactly the structure of our case categories, where morphisms
represent grammatical relations and composition encodes transitivity.
\end{frame}

\begin{frame}{Case Roles as Functional Typing for Agent Systems}
\protect\phantomsection\label{case-roles-as-functional-typing-for-agent-systems}
The eight-case framework of CEREBRUM {[}@friedman2024cerebrum{]} maps
directly onto the functional roles that components play in multi-agent
architectures:

{\def\LTcaptype{none} % do not increment counter
\begin{longtable}[]{@{}
  >{\raggedright\arraybackslash}p{(\linewidth - 4\tabcolsep) * \real{0.3333}}
  >{\raggedright\arraybackslash}p{(\linewidth - 4\tabcolsep) * \real{0.3333}}
  >{\raggedright\arraybackslash}p{(\linewidth - 4\tabcolsep) * \real{0.3333}}@{}}
\toprule\noalign{}
\begin{minipage}[b]{\linewidth}\raggedright
Case Role
\end{minipage} & \begin{minipage}[b]{\linewidth}\raggedright
Agent System Analogue
\end{minipage} & \begin{minipage}[b]{\linewidth}\raggedright
Protocol Function
\end{minipage} \\
\midrule\noalign{}
\endhead
NOM (Agent) & Active requester & Initiator of action policies \\
ACC (Patient) & Passive responder & Target of directed operations \\
GEN (Source) & Context provider & Supplier of prior information \\
DAT (Recipient) & Designated receiver & Endpoint for information
transfer \\
INS (Instrument) & Tool / API & Means of state transformation \\
LOC (Context) & Environment state & Boundary conditions and Markov
blanket \\
ABL (Origin) & Data source & Causal provenance of inputs \\
VOC (Addressee) & Routing target & Pragmatic addressing in broadcast
networks \\
\bottomrule\noalign{}
\end{longtable}
}

This mapping is not merely analogical. In the MCP protocol, a tool call
has precisely the structure of a case-marked clause: an \emph{agent}
(NOM) invokes a \emph{tool} (INS) on a \emph{resource} (ACC) to deliver
a \emph{result} to a \emph{recipient} (DAT). The case structure ensures
that every interaction is relationally typed, preventing the category
errors (sending a response where a tool call is expected) that plague
unstructured message-passing systems.
\end{frame}

\begin{frame}{Categorical Deep Learning: An Algebraic Theory of
Architectures}
\protect\phantomsection\label{categorical-deep-learning-an-algebraic-theory-of-architectures}
Gavranović {[}-@gavranovic2024thesis{]} develops ``a novel mathematical
foundation for deep learning based on the language of category theory,''
showing that neural network architectures---feedforward, convolutional,
recurrent, attention---can all be subsumed under a single categorical
framework using parameterized optics and lenses. This ``Categorical Deep
Learning'' approach {[}@shiebler2024categorical{]} demonstrates that
categorical structure is ``an algebraic theory of all architectures,''
unifying design patterns that appear unrelated when viewed through
conventional linear algebra.

For case theory in AI, this result is significant: it means the same
categorical language used to describe linguistic case assignment also
describes the internal computation of the neural networks that process
language. A transformer's attention mechanism, viewed through
Gavranović's lens, performs a kind of \emph{categorical case
assignment}---each attention head selects which input tokens play which
functional roles in the computation, with attention weights serving as
the enriched hom-values of \autoref{sec:enriched-categories}.
\end{frame}

\begin{frame}{String Diagrams as a Foundation for Interpretable AI}
\protect\phantomsection\label{string-diagrams-as-a-foundation-for-interpretable-ai}
The lambeq library {[}@lorenz2023lambeq{]} demonstrates that string
diagrams provide a practical interface between linguistic structure and
machine learning. As an ``efficient high-level Python library for
Quantum NLP,'' lambeq translates sentences into string diagrams,
converts diagrams into parameterized circuits (quantum or classical),
and trains parameters end-to-end on NLP tasks. The diagrammatic
representation serves dual purposes: it is both the mathematical
specification of the model \emph{and} a human-readable explanation of
what the model computes.

This interpretability property is crucial for AI safety and alignment.
Where transformer architectures produce opaque attention patterns, a
DisCoCat model deployed via string diagrams produces derivation trees
with explicit compositional semantics---every box and wire has a
linguistic interpretation. Extending this to agent communication, a
categorical protocol would produce not just working message exchanges
but \emph{interpretable} interaction diagrams where each step's
relational role is transparent.
\end{frame}

\begin{frame}{DisCoCirc and Multi-Turn Agent Dialogue Protocols}
\protect\phantomsection\label{discocirc-and-multi-turn-agent-dialogue-protocols}
The DisCoCirc framework {[}@defelice2022discocirc{]} extends
compositional semantics from single sentences to multi-sentence
discourse by introducing \emph{state wires} that persist across sentence
boundaries. This architecture maps directly onto multi-turn agent
dialogues:

\begin{itemize}
\tightlist
\item
  \textbf{Entity wires} correspond to persistent agent identities across
  communication rounds
\item
  \textbf{State updates} correspond to belief revisions triggered by
  incoming messages
\item
  \textbf{Coreference resolution} corresponds to entity tracking across
  protocol sessions
\item
  \textbf{Discourse coherence} corresponds to protocol correctness
  constraints
\end{itemize}

In a multi-agent system, a DisCoCirc-style protocol would track the
evolving states of all participating agents as persistent wires in a
circuit diagram, with each message exchange represented as a box that
transforms the relevant wires. Protocol correctness reduces to a
categorical property: the circuit must type-check, meaning all wire
types must match at connection points---precisely the condition that
case marking enforces in natural language.
\end{frame}

\begin{frame}{Compositional Game Theory and Categorical Multi-Agent
Equilibria}
\protect\phantomsection\label{compositional-game-theory-and-categorical-multi-agent-equilibria}
The ``parametrised optics'' framework developed within categorical
cybernetics ``provides a general-purpose foundation for the study of
controlled processes'' applicable to compositional game theory as a
multi-agent framework. In this setting, agents are modeled as lenses (or
optics) that observe the environment through one channel and act through
another, with the overall system behavior emerging from the composition
of individual agent behaviors.

This connects to our enriched case framework
(\autoref{sec:enriched-categories}): the precision weights on morphisms
correspond to the utility parameters of game-theoretic agents, and the
composition inequality
\(\mathcal{C}(A,C) \geq \mathcal{C}(A,B) \cdot \mathcal{C}(B,C)\)
corresponds to the sub-optimality of indirect communication chains.
Equilibria in the multi-agent game correspond to fixed points of the
enriched functor, where no agent can improve its utility by changing its
case-role assignment.
\end{frame}

\begin{frame}{Double Categorical Systems Theory for Explainable
Autonomous AI}
\protect\phantomsection\label{double-categorical-systems-theory-for-explainable-autonomous-ai}
Recent work on Double Categorical Systems Theory (DCST) aims to
``utilise Double Categorical Systems Theory as a mathematical framework
to facilitate collaboration in co-designing an explainable and auditable
model of an autonomous AI system's deployment environment.'' DCST
extends ordinary category theory with a second dimension of morphisms
(2-morphisms), allowing simultaneous modeling of:

\begin{enumerate}
\tightlist
\item
  \textbf{Horizontal composition}: Sequential chaining of agent
  interactions (morphism composition)
\item
  \textbf{Vertical composition}: Hierarchical nesting of subsystems
  within larger systems (2-morphism composition)
\end{enumerate}

For case theory, the double-categorical extension allows us to model
both the \emph{within-sentence} case structure (horizontal) and the
\emph{discourse-level} case structure (vertical) within a single
algebraic framework---precisely the unification that DisCoCirc achieves
pragmatically.
\end{frame}

\begin{frame}{Toward Categorical Communication Standards for AI Agents}
\protect\phantomsection\label{toward-categorical-communication-standards-for-ai-agents}
The synthesis of these developments suggests a research program:
developing \textbf{categorical communication protocols} that combine the
engineering robustness of existing standards (A2A, MCP, ACP) with the
compositional semantics of categorical linguistics. Such a protocol
would:

\begin{enumerate}
\tightlist
\item
  \textbf{Type-check interactions}: Every message exchange would be
  relationally typed by case roles, preventing structural communication
  errors at the protocol level
\item
  \textbf{Compose transparently}: Multi-step interactions would compose
  algebraically, with diagrammatic representations providing
  interpretable audit trails
\item
  \textbf{Transfer across implementations}: Topos-theoretic bridges
  (\autoref{sec:topos-theory}) would ensure that a protocol verified in
  one categorical formalization automatically satisfies correctness
  conditions in any Morita-equivalent formulation
\item
  \textbf{Scale via enrichment}: Distributional proximity measures
  (\autoref{sec:enriched-categories}) would enable graceful degradation
  under uncertainty, with enriched weights encoding confidence in
  message content
\item
  \textbf{Ground in cognitive architecture}: The active inference
  foundation (\autoref{sec:cognitive-integration}) would ensure that
  artificial agents communicate using the same relational structure that
  evolution has optimized for biological cognition
\end{enumerate}
\end{frame}

\begin{frame}{Quantum Cryptographic Communications and Semantic
Security}
\protect\phantomsection\label{quantum-cryptographic-communications-and-semantic-security}
The quantum topological framework of
\autoref{sec:quantum-active-inference} opens a natural connection to
quantum cryptographic communications, with implications for both the
security and the semantic integrity of case-structured information
transfer.

\begin{block}{Quantum Key Distribution for Relational Semantics}
\protect\phantomsection\label{quantum-key-distribution-for-relational-semantics}
Quantum key distribution (QKD) protocols provide information-theoretic
security guarantees that classical cryptography cannot achieve
{[}@pirandola2020qkd{]}. When case-marked relational structures are
communicated between agents---whether human or artificial---QKD ensures
that the \emph{relational semantics} of a message (who does what to
whom) cannot be intercepted or altered without detection. This is
particularly relevant for high-stakes domains where case assignment
carries legal or medical significance: a tampered case frame could
change an agent's interpretation of responsibility, causation, or
obligation.

The sheaf-theoretic framework of Khatri et al. {[}-@khatri2025quantum{]}
further shows that entanglement-assisted semantic channels strictly
exceed classical semantic capacity:

\begin{quote}
``We derive semantic channel capacity when sender and receiver share
prior entanglement, proving it strictly exceeds classical capacity.''
{[}@khatri2025quantum{]}
\end{quote}

This means that quantum-secured channels not only protect relational
content but also enable \emph{more} relational content to be
communicated per channel use---a genuine quantum advantage for semantic
communication.
\end{block}

\begin{block}{Semantic Cryptography: Encrypting Meaning Structures}
\protect\phantomsection\label{semantic-cryptography-encrypting-meaning-structures}
Beyond bit-level QKD, the categorical framework suggests a notion of
\emph{semantic cryptography}: encrypting not just the symbols of a
message but its compositional meaning structure. In a DisCoCat
framework, a sentence's meaning is a morphism in a compact closed
category---a multilinear map from word spaces to sentence space.
Semantic encryption would operate on this categorical level:

\begin{enumerate}
\tightlist
\item
  \textbf{Functorial encryption}: Applying a secret functor
  \(F\colon \mathbf{C} \to \mathbf{D}\) that maps the plaintext case
  category into a ciphertext category, preserving compositional
  structure but rendering individual meanings unintelligible without the
  inverse functor.
\item
  \textbf{Diagram obfuscation}: Applying ZX-style rewrites that change
  the internal topology of a DisCoCat derivation while preserving its
  overall semantics---creating multiple equivalent ``ciphertexts'' for
  the same semantic ``plaintext,'' each with a different diagrammatic
  structure.
\item
  \textbf{Enriched weight masking}: In an enriched case category, the
  hom-values (distributional weights) can be encrypted independently of
  the categorical structure, allowing transmission of relational
  topology without revealing the distributional content.
\end{enumerate}

These operations extend the cryptographic primitives beyond QKD into
genuinely compositional territory {[}@broadbent2016qcrypto{]}, where the
mathematical structure of categorical semantics provides the algebraic
substrate for security proofs.
\end{block}

\begin{block}{Cognitive Security: Protecting Relational Reasoning}
\protect\phantomsection\label{cognitive-security-protecting-relational-reasoning}
The cognitive integration of \autoref{sec:cognitive-integration} raises
a complementary concern: \emph{cognitive security}---ensuring that an
active inference agent's generative model of relational structure is not
adversarially corrupted. In the predictive processing framework, an
agent's case-assignment system is a generative model that minimizes
variational free energy. Adversarial manipulation of this system could:

\begin{itemize}
\tightlist
\item
  \textbf{Inject false case frames}: Leading an agent to misidentify the
  agent, patient, or instrument of an action---a form of semantic
  adversarial attack
\item
  \textbf{Exploit precision weighting}: Artificially inflating the
  precision of misleading sensory evidence, causing the agent to update
  its case assignments toward adversarially chosen interpretations
\item
  \textbf{Corrupt topos-theoretic transfer}: If an agent uses Morita
  equivalence to transfer case-theoretic results between frameworks,
  corrupting one framework could propagate errors across all equivalent
  formulations
\end{itemize}

Defending against these attacks requires \emph{quantum-secured cognitive
integrity}: using quantum authentication and tamper-detection protocols
to ensure that the generative model's parameters---the objects,
morphisms, and enriched weights of the case category---have not been
adversarially modified. The topological robustness of TQNNs
(\autoref{sec:quantum-active-inference}) provides a natural resilience
mechanism: topological invariants are robust to continuous deformation,
so small perturbations to the generative model's parameters cannot
change its topological class, providing inherent resistance to
gradient-based adversarial attacks.
\end{block}
\end{frame}

\end{document}
